\chapter{Технологии, репозиториуми и почетен систем}
\label{chap:technology}

\section{Развојна средина}

Развојниот host е WSL2 со Ubuntu 24.04. ROS Noetic, Gazebo Classic и catkin се
извршуваат во Docker image базиран на Ubuntu 20.04. Ваквата комбинација го
задржува поддржаниот Noetic userspace, а истовремено овозможува современ host и
WSLg приказ. Главниот container го монтира Lite3 workspace-от read--write и
EmotionBot read-only. Live OpenAI патот е посебен disposable Python 3.12
container; offline режимот и сите редовни тестови не бараат API key, интернет,
GUI, микрофон или физички робот.

Во \cref{tab:technologies} се сумирани главните технологии. ROS Noetic и
Gazebo се користат според нивните официјални интерфејси; MotionSDK ограничувањата и безбедносните
предупредувања се преземени од официјалниот репозиториум и локалните manuals.

\begin{table}[htbp]
\centering
\caption{Користени технологии и нивната улога}
\label{tab:technologies}
\small
\begin{tabularx}{\textwidth}{L{0.22\textwidth}YY}
\toprule
\textbf{Технологија} & \textbf{Улога} & \textbf{Граница/напомена} \\
\midrule
DEEP Robotics Lite3 & Физичка и симулациска четириножна платформа & Физичкиот robot path не се активира од симулациски target. \\
ROS Noetic & Nodes, topics, services, параметри и orchestration & Работи во Ubuntu 20.04 Docker, не се инсталира на WSL host. \\
Gazebo Classic & Динамичка симулација, joint/contact набљудување и визуелна проверка & Резултатите не се физичко commissioning evidence. \\
Docker/WSL2 & Репродуцибилна Ubuntu 20.04 средина врз Ubuntu 24.04 WSL2 & GUI се прикажува преку WSLg. \\
C++14/catkin & Lite3 controller, траектории и ROS build & Се задржува компатибилноста со upstream. \\
Python 3.8 & ROS adapters, mapper, safety bridge и тестови & Noetic средината не го добива целосниот EmotionBot dependency set. \\
Python 3.12 sidecar & OpenAI SDK и Responses streaming & Loopback-only; API key не влегува во ROS workspace. \\
EmotionBot и EmotionEngine & Appraisal, personality, memory, valence, arousal и category & GPL-3.0; роботска кореографија не се сместува во овој репозиториум. \\
MotionSDK & Единствен физички joint-command owner & Бара свеж feedback, state gates, STOP и детерминистички release. \\
Git submodules & Pinned, репродуцибилни checkouts & Nested commit-и се објавуваат пред wrapper gitlinks. \\
\bottomrule
\end{tabularx}
\end{table}


\section{Git структура и референтен snapshot}

Wrapper-репозиториумот pin-ира три submodules. Одржуваните изворни
репозиториуми се EmotionBot и Lite3\_VMC fork-от; \code{lite3\_vmc\_upstream}
е чиста споредбена копија и во неа не се развива. Структурата е прикажана на
\cref{fig:repo-structure}. Јавните URL-адреси, гранките и SHA-вредностите биле
проверени на 20 септември 2026 година.

\begin{figure}[htbp]
\centering
\resizebox{\textwidth}{!}{\begin{tikzpicture}[node distance=11mm and 12mm]
  \node[box,text width=45mm] (wrapper) {
    \textbf{lite3-ros-sim}\\
    wrapper, Docker, docs и tickets\\
    \code{c41f26e}};

  \node[process,below left=of wrapper,text width=39mm] (emotion) {
    \textbf{emotion-bot}\\
    headless EmotionEngine\\
    \code{f267d28}};
  \node[process,below=of wrapper,text width=39mm] (fork) {
    \textbf{Lite3\_VMC fork}\\
    controller + emotion\_bot\_ros\\
    \code{785a384}};
  \node[process,below right=of wrapper,text width=39mm] (upstream) {
    \textbf{Lite3\_VMC upstream}\\
    clean reference\\
    \code{724aa54}};

  \draw[flow] (wrapper.south west) -- node[left,font=\scriptsize]{gitlink} (emotion.north);
  \draw[flow] (wrapper.south) -- node[right,font=\scriptsize]{gitlink} (fork.north);
  \draw[flow] (wrapper.south east) -- node[right,font=\scriptsize]{gitlink} (upstream.north);
  \coordinate (compareleft) at ($(fork.south)+(0,-5mm)$);
  \coordinate (compareright) at (compareleft -| upstream.south);
  \draw[dashedflow] (fork.south) -- (compareleft) -- (compareright) -- (upstream.south);
  \node[below=1mm,font=\scriptsize,fill=white,inner sep=1.5pt]
    at ($(compareleft)!0.5!(compareright)$) {само споредба; без развој};
\end{tikzpicture}
}
\caption{Wrapper и pinned submodule структура на референтниот snapshot.}
\label{fig:repo-structure}
\end{figure}

Wrapper-от е commit
\code{c41f26ef4414e5c9a4748f3c88aee33820526c01}. Неговиот EmotionBot gitlink
покажува кон \code{f267d28552bbf46f30a263bb291e71ab91a2b91b}, а активниот
Lite3\_VMC fork кон \code{785a3846ec2c44a0380618490f837786b3250f76}. Upstream
референцата е \code{724aa54dee0cc374bd5417c86d649fb08d4d309a}. Овие commit-и
се важни затоа што датиран verification record не докажува дека произволен
подоцнежен checkout има исти резултати.

\section{Почетниот Lite3\_VMC систем}

Оригиналниот DEEP Robotics Lite3\_VMC обезбедува Gazebo модел, joint
controllers, model spawner, keyboard input и симулациски executable. Baseline
проверката утврдила:

\begin{itemize}
  \item 13 активни контролери: 12 joint effort controllers и еден
  joint-state controller;
  \item 12 joints на \code{/lite3\_gazebo/joint\_states};
  \item keyboard publisher на \code{sensor\_msgs/Joy} преку \code{/joy}; и
  \item симулацискиот executable го инстанцира Joy receiver-от, но не го
  инстанцира постојниот \code{/cmd\_vel} receiver и не го применува
  \code{speed\_update\_mode}.
\end{itemize}

Последната точка е суштинска: присуството на source class не значи дека
runtime path е имплементиран. Поради тоа новата интеграција го ремапира
единствениот контролен влез на \code{/emotion\_bot/joy\_out}; во трудот не се
тврди дека \code{/cmd\_vel} е работен симулациски пат.

Upstream spawner-от е чувствителен на standard input: првиот Enter ги стартува
контролерите, а вториот ги запира и го отстранува моделот. EOF може веднаш да
предизвика cleanup. Интегрираниот launch затоа користи service-based controller
loader и supervising runner, без да го менува upstream baseline workflow.

\section{EmotionBot и лиценцна граница}

EmotionBot е GPL-3.0 проект. Во интеграцијата не се копира неговиот emotional
core, ниту интерактивниот \code{main.py} се автоматизира преку stdin. Наместо
тоа, во EmotionBot е додаден повторно употреблив headless API
\code{emotional\_core.engine.EmotionEngine}. Тој управува со appraisal,
personality, memory, valence, arousal, categorical state и deterministic
fallback. Import-от не бара plotting, микрофон, network, модел или API key.

ROS contract-ите, conversation ordering, mapping и симулациска заштита се во
новиот \code{emotion\_bot\_ros} package во Lite3\_VMC fork-от. Физичкиот
hardware workspace, MotionSDK runner и ticket evidence се во wrapper-от. Со
ова се избегнува погрешно припишување: upstream Lite3 controller и EmotionBot
domain logic остануваат дела со сопствена историја и лиценца, додека
интеграциските слоеви се проектни измени.

\section{Барања и ограничувања}

Главните функционални барања се девет категории, live и offline разговор,
верзионирани состојби, визуелно различни Gazebo профили, exact-neutral
транзиции и inspectable status. Нефункционалните барања се детерминизам на
тестовите, bounded latency/size, no-secret handling, fail-safe zeroing,
ексклузивна hardware ownership и репродуцибилен build.

Ограничувањата се намерно строги. Normal chat не задава planar locomotion или
yaw; Gazebo hop/stomp не се пренесуваат на hardware; physical airborne hop,
torque strike и external wrench не се дозволени; vendor \code{lite\_cog},
\code{qnx2ros} и \code{jy\_exe} не се менуваат. Физичката работа не е нормален
дел од симулацискиот workflow и бара одделна авторизација, operator и fresh
telemetry.
